% RESEARCH_DRAFT_SHA256: f8c6f99b38d262a1e69003c04007053f33f750bd9238d6ae10ac25d66d22a123
% Layout and source-reference wrappers expanded for the web renderer.
\documentclass[11pt]{article}
\usepackage[margin=25mm]{geometry}
\usepackage{amsthm}
\usepackage{fontspec,unicode-math}
\setmainfont{Latin Modern Roman}
\setmathfont{Latin Modern Math}
\usepackage{amsmath,amssymb,mathtools}
\usepackage{xurl,hyperref}
\hypersetup{colorlinks=true,urlcolor=blue,linkcolor=blue}
\setcounter{secnumdepth}{0}
\setlength{\parindent}{0pt}
\setlength{\parskip}{5pt}
\setlength{\emergencystretch}{3em}
\newcommand{\sref}[2]{\hyperlink{source:#1:#2}{[S#2]}}
\newcommand{\cataloguescope}[1]{\paragraph{Source question.}#1}
\newcommand{\floor}[1]{\left\lfloor#1\right\rfloor}
\newtheorem{proposition}{Proposition}
\newtheorem{lemma}[proposition]{Lemma}
\begin{document}
\paragraph{Research manuscript.} OpenAI \texttt{gpt-6-astra}, reasoning \texttt{ultra}; 11 October 2026.
\paragraph{Preserved source body.} The complete catalogue mathematical text and sources follow with layout and source links expanded, including source qualifications and its historical preparation notes. New research begins at the section headed Research attempt.

% UnsolvedMath ID 30005832; source catalogue code OWR-14298166-010
\section{Fractional Coefficient Savings in Algebraic Proof Systems}\label{30005832}
{\small UnsolvedMath ID 30005832 · Computer Science}
\subsection{Definitions and mathematical statement}

\paragraph{Shared notation.}Standard number systems are denoted by $\mathbb N,\mathbb Z,\mathbb Q,\mathbb R,\mathbb C$; index ranges, computational models and specialized notation follow the cited problem source and the definitions below.


For equations $f_i=0$ with rational coefficients, a Nullstellensatz refutation has an identity $\sum_i g_i f_i=1$ with polynomial coefficients $g_i$ over the permitted coefficient field. For Boolean problems the equations $x_j^2-x_j=0$ are included in the source encoding. Sherali--Adams certificates combine the input inequalities with products of Boolean literals under a degree limit, using the source normalization. The total coefficient size counts the binary encoding of all rational numerators and denominators (or integers in the integral restriction), along with their occurrence in the certificate. Fractional coefficients are nonintegral rational coefficients.
\paragraph{Statement recorded in the supplied source.}
Are there natural examples in which fractional coefficients greatly reduce the total coefficient size required for Nullstellensatz or Sherali–Adams proofs? \href{https://doi.org/10.4171/owr/2024/15}{[S2]}
\paragraph{Scope qualification.}
The normalization and integral comparison must be fixed: clearing denominators changes the right-hand constant in an identity normalized to 1. Certificate degree, monomial support and coefficient bit size are distinct measures; the question specifically asks about coefficient savings.
\paragraph{Source status.}
The supplied catalogue labels this record partially solved. Its source-status metadata was checked against the supplied snapshot on 2026-10-11. This import does not independently establish current open status or mathematical correctness. The catalogue literature assessment is dated 2026-08-22.
\subsection{Short English statement}
Can natural constraint families have much shorter algebraic certificates when noninteger rational coefficients are permitted?
\subsection{Sources}
\begin{itemize}
\item[S1] Ulam AI and the UnsolvedMath Contributors, UnsolvedMath: A Curated Collection of Open Mathematics Problems, record 30005832 (OWR-14298166-010). Catalogue curation and metadata: CC BY 4.0. The supplied source question is retained; definitions and the English summary are editorial additions.
\url{https://huggingface.co/datasets/ulamai/UnsolvedMath}.
\item[S2] Proof Complexity and Beyond (2024), Oberwolfach Reports, DOI 10.4171/owr/2024/15, pp. 871--934
\url{https://doi.org/10.4171/owr/2024/15}.
\item[S3] Potechin and Zhang, Bounds on Total Coefficient Size
\url{https://doi.org/10.4230/LIPIcs.ICALP.2024.117}.
\item[S4] Potechin and Zhang, full version
\url{https://arxiv.org/abs/2205.03577}.
\textit{Underlying Oberwolfach Reports material retains its source-specific licence, including CC BY-SA 4.0 where applicable; see THIRD\_PARTY\_NOTICES.md.}
\end{itemize}
\paragraph{Preparation and review.}
Original question attributed to the cited source and UnsolvedMath contributors. Codex (GPT-6; exact serving revision not exposed) assisted with definitions, English summaries, source inspection and formatting. No new human mathematical authorship or independent proof verification is asserted. Canonical ID 30005832 and source code OWR-14298166-010 checked against the complete supplied registry and existing repository statements; known flagged repetitions and obvious identical conjectures excluded. Full mathematical-equivalence/literature review has not been performed. Statement preparation and final review followed contribution prompts C and D; automated checks verify structure and identity, not open status.
\label{end:30005832}

\section{Research attempt}
\subsection{Literature, source context and scope}
The full original catalogue body and source qualifications are preserved above. The relevant primary workshop passages were read in the cached publisher report. Historical claims and subsequent attributed results are distinguished from the mathematical work below. A restricted example, standard lemma or known theorem is not counted as new progress on the entire original target. Missing assumptions and repeated or rhetorical extracts are not silently turned into different solved problems.

\subsection{Normalization separates coefficient magnitude from bit length}
The preserved catalogue describes coefficient bit complexity. The primary report on pp.926--927 instead defines total coefficient size by the sum of absolute coefficient values. Potechin--Zhang's primary abstract also concerns this magnitude measure. These are different questions; the original text above is preserved rather than silently changed.
We test the tempting inference that clearing rational denominators preserves small normalized certificates. For an integer $M\ge2$, take the deliberately elementary inconsistent constant equations $M=0$ and $M+1=0$, over characteristic zero. A Nullstellensatz certificate with constant multipliers is
\[
Ma+(M+1)b=1.
\]
This is a normalization experiment, not a proposed natural hard Boolean family.
\begin{proposition}
The least value of $|a|+|b|$ is $1/(M+1)$ for rational multipliers and $2$ for integer multipliers. Allowing polynomial multipliers, while counting all their coefficient magnitudes, does not lower either optimum.
\end{proposition}
\begin{proof}
For rationals, $1\le M|a|+(M+1)|b|\le(M+1)(|a|+|b|)$, and $a=0,b=1/(M+1)$ attains equality. For integers, $a=-1,b=1$ gives value two. An integer pair of magnitude sum less than two is zero or has a single coefficient equal to $1$ or $-1$, which cannot produce one because $M\ge2$. For polynomial multipliers, taking constant terms in the polynomial identity yields exactly the displayed scalar equation; total coefficient magnitude dominates that of the constant terms. Constant multipliers attain the preceding bounds.
\end{proof}
The magnitude ratio is $2(M+1)$, whereas the rational certificate writes a denominator requiring $\Theta(\log M)$ bits and the integer certificate uses two unit coefficients. Thus even the direction of a size comparison can change with the measure. Clearing the denominator of the rational identity gives right-hand side $M+1$, not the required normalized right-hand side one. Dividing again leaves the integer setting. This is the precise failed step in that proposed normalization argument.
\subsection{Verification and remaining gaps}
The optima were checked by bounded integer enumeration and exact rational arithmetic for several $M$, including $M=2^t$. The proof, not the enumeration, establishes global optimality. The input constants themselves also have encoding cost $\Theta(\log M)$; that cost has not disappeared. The example gives neither a bit-complexity lower bound for the catalogue question nor a separation for natural Boolean or Sherali--Adams families. Such a result requires the correct measure, an agreed encoding, and a family for which the auxiliary normalization effect is meaningful.

\section{References and source attribution}
Primary workshop report: \url{https://doi.org/10.4171/owr/2024/15}. Relevant passages read: 2024/15 pp.926--927, magnitude definition and rational-versus-integral coefficient questions.

\url{https://arxiv.org/abs/2205.03577}. Read primary Potechin--Zhang abstract identifying exponential total-coefficient-size results; this magnitude statement is not relabeled as a bit-complexity theorem.

All retrievals were checked on 11 October 2026. Original question attribution remains in the preserved Sources section. External theorems are attributed and not claimed as new discoveries or fully verified merely by citation.


\paragraph{Integration review.} Codex (GPT-6; exact serving revision not exposed) reviewed the written arguments and expanded layout and source-reference wrappers. Repository metadata and the production web renderer were checked using Python 3.13 and JavaScript with MathJax 3.2.2. This model-assisted review is not independent human proof verification; no PDF compilation or formal proof checking is asserted.

\section{Completion Estimate}
\noindent\textbf{COMPLETION: 0\%}
No substantiated new progress toward the entire original target is claimed. The worked results and failed route are delimited above; a known result, illustrative model, or auxiliary obstruction is not counted as original-target completion.

\section{Final status and attribution}
\textbf{Attempt status: unresolved.} This describes this research attempt, not a new ruling on the source catalogue's status or on external literature claims.
The mathematics and this self-audit were generated by OpenAI \texttt{gpt-6-astra}, reasoning \texttt{ultra}, on 11 October 2026. Codex (GPT-6; exact serving revision not exposed) orchestrated the batch, read and reviewed its written mathematical arguments, and contributed formatting/repository review. No human mathematical authorship, independent human proof checking, or formal verification is claimed. The historical preparation paragraphs in the reproduced catalogue body concern that body only.

Contribution prompt A was applied to the complete definition and supplied source record, followed by the adversarial checks of prompt B. Local reproduction files are \texttt{batch025/checks.py}, \texttt{computations.json}, and the per-ID audit JSON. They were not public links at preparation time. Numerical checks are finite and do not prove unrestricted results. No LaTeX engine was found locally; compilation was not performed by this research agent.

\end{document}
